\begin{textsample}{POS Dim 4 – 2017\_c6 – Score 10.00 – 2017\_c6/enade\_2017\_art\_vis\_8100b5f9fc.txt}  \label{ex:f4_pos_040}
CAPÍTULO 6 ANÁLISE TÉCNICA DA PROVA Este capítulo tem por objetivo apresentar o desempenho dos estudantes concluintes de Artes Visuais ( Licenciatura ) no Enade/2017.
Para isso, foram calculadas as estatísticas básicas da prova em seu todo ( seção 6.1.1 ), bem como as estatísticas relacionadas aos Componentes de Formação Geral ( seção 6.1.2 ) e de Conhecimento Específico da Área ( seção 6.1.3 ).
Dadas as suas características, foram analisadas, em separado, as questões objetivas ( seção 6.2 ) e as questões discursivas ( seção 6.3 ).
Para as questões objetivas, foram disponibilizados os índices de facilidade e de discriminação ponto-bisserial, também em separado, para os Componentes de Formação Geral ( seção 6.2.1 ) e de Conhecimento Específico ( seção 6.2.2 ).
De cada componente, uma das questões foi escolhida para exemplificar a análise gráfica, relacionando as alternativas escolhidas pelos estudantes ( inclusive o gabarito ) com o número de acertos no componente.
No Anexo I, apresenta -se a íntegra da análise gráfica para todas as questões objetivas.
Para cada uma das questões discursivas, os conteúdos dos tipos mais comuns de respostas dos estudantes são apresentados e comparados com o padrão de respostas esperado ( ver Anexo VIII com o padrão de respostas ).
Tomando -se como base as duas questões discursivas do Componente de Formação Geral, na seção 6.3.1.6 apresentam -se comentários sobre a correção das respostas com respeito à Língua Portuguesa.
Nas tabelas são apresentadas as seguintes estatísticas das notas¹⁹ : média do desempenho na prova, erro padrão da média, desvio padrão, nota mínima, mediana e nota máxima para cada um de seus componentes.
Tais estatísticas contemplam o total de estudantes concluintes da Área de Artes Visuais ( Licenciatura ) inscritos e presentes à prova do Enade/2017, tendo em vista agregações, ou por Grandes Regiões e o país como um todo, ou por Categoria Administrativa, Organização Acadêmica da IES e Modalidade de Ensino.
Em relação aos gráficos de distribuição de notas, o intervalo considerado foi de 10 unidades, aberto à esquerda e fechado à direita, com exceção do primeiro intervalo, [ 0 ; 10 ], fechado em ambos os extremos.
Para os gráficos de distribuição das notas das questões discursivas, foram consideradas mais duas categorias : questão em branco²⁰ e nota zero. ¹⁹ Uma definição dessas estatísticas pode ser encontrada no Glossário. ²⁰ Nesse grupo estão incluídas também as respostas classificadas como nulas ou desconsideradas. 6.1 ESTATÍSTICAS BÁSICAS DA PROVA Esta seção apresenta estatísticas selecionadas e histogramas da nota geral ( 6.1.1 ) e de cada componente : Formação Geral ( 6.1.2 ) e Conhecimento Específico ( 6.1.3 ).
São também apresentadas estatísticas selecionadas de subpopulações caracterizadas por Grande Região, Categoria Administrativa, Organização Acadêmica e Modalidade de Ensino. 6.1.1 Estatísticas Básicas Gerais Na Tabela 6.1, apresentam -se as Estatísticas Básicas da prova, por Grande Região, dos estudantes concluintes de Artes Visuais ( Licenciatura ).
A Média das notas da prova como um todo ( nas seções seguintes serão analisados os Componentes de Formação Geral e de Conhecimento Específico ) foi 51,4, sendo que os alunos da região Sul obtiveram a Média mais baixa ( 48,7 ), e os da região Centro-Oeste obtiveram a Média mais alta ( 54,7 ).
As demais médias foram : 50,3 na região Norte, 53,9 na região Nordeste e 52,7 na região Sudeste.
O Desvio padrão para o Brasil como um todo foi 15,2, sendo o maior Desvio padrão encontrado na região Sudeste ( 15,4 ), e o menor, na região Norte ( 14,4 ), indicando uma dispersão um pouco menor das notas dessa última região.
A região que obteve a maior nota Máxima foi a Sudeste ( 90,1 ), ao passo que a região que atingiu a menor nota Máxima foi a Centro-Oeste ( 81,7 ).
A Mediana do Brasil como um todo foi 52,1, sendo a maior Mediana obtida na região Nordeste ( 56,3 ), e a menor obtida na Sul ( 48,6 ).
A nota Mínima foi zero na região Sudeste, 7,9 na região Sul, 13,8 na região Centro-Oeste, 16,3 na região Norte e 19,6 na região Nordeste.
Considerando -se as notas segundo Grande Região, observa -se que existe diferença estatisticamente significativa ao nível de 95 \% entre a menor Média, obtida na região Sul ( 48,7 ), e as médias das regiões Nordeste, Sudeste e Centro-Oeste²¹. ²¹ Todas as comparações deste capítulo utilizam os intervalos de 95 \%.
Os erros-padrão da média que possibilitam os testes estão disponíveis nas tabelas.
O comportamento das notas dos estudantes de todo o Brasil pode ser observado no Gráfico 6.1, que apresenta um histograma com a distribuição dessas notas.
Essa é uma distribuição unimodal com a moda no intervalo ( 50 ; 60 ], seguido pelo intervalo ( 40 ; 50 ].
Na Tabela 6.2, apresentam -se informações referentes à Média da nota final desagregadas por Categoria Administrativa, por Organização Acadêmica e por Modalidade de Ensino.
A Média das notas da prova como um todo foi 51,4.
Em relação à Categoria Administrativa, os estudantes das IES Públicas obtiveram Média mais alta ( 55,9 ), e os das IES Privadas obtiveram Média mais baixa ( 49,4 ), que a Média nacional.
Observa -se que existe diferença estatisticamente significativa entre as médias das notas das IES Públicas e Privadas.
A diferença entre as médias das regiões Centro-Oeste e Sul ( 6,0 ), a maior e a menor Média, é inferior à diferença entre IES Públicas e Privadas ( 6,5 ), caracterizando uma maior diversidade administrativa do que regional.
O Desvio padrão para as IES Públicas ( 15,0 ) foi inferior ao do Brasil como um todo ( 15,2 ), indicando uma dispersão um pouco menor das notas nessa Categoria Administrativa.
No tocante à Organização Acadêmica, as Universidades e os CEFET/IFET obtiveram Média mais alta que a nacional ( 53,3 e 55,2, respectivamente ).
A Média dos Centros Universitários foi a menor : 48,8.
Constata -se que existe diferença estatisticamente significativa ao nível de 95 \% entre a Média mais baixa, 48,8 dos Centro Universitários, e a dos outros tipos de Organização Acadêmica.
A Média da Modalidade Educação a Distância ( 46,8 ) foi inferior à Média da Educação Presencial ( 55,3 ), e a diferença entre elas é estatisticamente significativa ao nível de 95 \%. 6.1.2 Estatísticas Básicas no Componente de Formação Geral Na Tabela 6.3, são apresentadas as Estatísticas Básicas em relação ao componente da prova que avalia a Formação Geral dos estudantes concluintes.
Os alunos de todo o Brasil obtiveram desempenho médio de 49,0.
Quanto à variabilidade, o Desvio padrão das notas dos estudantes do Brasil, como um todo, foi 16,9.
A maior Média foi obtida na região Centro-Oeste ( 51,7 ), e a menor, na região Sul ( 47,8 ).
As demais médias foram : 48,9 na região Norte, 49,1 na região Nordeste e 49,6 na região Sudeste.
O maior Desvio padrão foi obtido na região Centro-Oeste ( 18,3 ), e o menor, na região Norte ( 15,2 ).
Os demais desvios padrões foram : 17,3 na região Nordeste, 17,2 na região Sudeste e 16,3 na região Sul.
A maior nota no Componente de Formação Geral da prova do Enade/2017 foi 94,2, obtida por, pelo menos, um aluno na região Sul.
A menor nota Máxima foi obtida na região Norte ( 90,1 ).
A Mediana do Brasil, como um todo, foi 49,9, sendo a menor Mediana encontrada na região Sul ( 48,4 ), e a maior, encontrada na região Centro-Oeste ( 52,3 ).
A nota Mínima nessa parte foi zero em quatro das cinco regiões, com exceção da região Norte, com nota Mínima igual a : 12,5.
Considerando -se as notas segundo Grande Região, observa -se que não existe diferença estatisticamente significativa entre as médias das notas do Componente de Formação Geral, obtidas nas diferentes regiões, a única exceção é entre as notas das regiões Sul e Sudeste.
No Gráfico 6.2, pode -se conferir a avaliação do desempenho dos estudantes do Componente de Formação Geral, a partir do histograma da distribuição das notas correspondentes.
A distribuição é unimodal, com moda em ( 50 ; 60 ], mesmo intervalo modal da distribuição de notas da prova, como um todo ( Gráfico 6.1 ).
Nota -se, ainda, que, no Gráfico 6.2, as notas apresentam dispersão maior do que a do Gráfico 6.1 ( distribuição das notas da prova ), confirmado pela comparação dos desvios padrões : 15,2 para a nota da prova como um todo e 16,9 para o Componente de Formação Geral.
Na Tabela 6.4 são apresentadas as informações referentes ao desempenho dos concluintes do Componente de Formação Geral, em diferentes agregações : Categoria Administrativa, Organização Acadêmica e Modalidade de Ensino.
Observa -se que existe diferença estatisticamente significativa entre as médias dos tipos de Categoria Administrativa : a maior Média obtida por estudantes de IES Públicas ( 52,6 ) e a menor, pelos de IES Privadas ( 47,4 ).
Nota -se que também não há diferença estatisticamente significativa entre as médias dos quatro tipos de Organização Acadêmica, valor maior para CEFET/IFET ( 51,4 ) e menor para Faculdades ( 46,8 ).
No que se refere à Modalidade de Ensino, a Educação Presencial apresentou maior Média do que a Educação a Distância ( médias 51,3 e 46,4, respectivamente ).
Tais médias apresentam diferença estatisticamente significativa. 6.1.3 Estatísticas Básicas do Componente de Conhecimento Específico Na Tabela 6.5, são apresentadas as Estatísticas Básicas referentes ao Componente de Conhecimento Específico da área de Artes Visuais ( Licenciatura ).
A Média do desempenho dos alunos do Brasil, como um todo, foi 52,2.
A maior Média foi obtida na região Centro-Oeste ( 55,7 ), e a menor, na região Sul ( 48,9 ).
As demais médias foram : 50,8 na região Norte, 55,5 na região Nordeste e 53,7 na região Sudeste.
Quanto à variabilidade das notas, o Desvio padrão do Brasil, como um todo, foi 16,8, o maior Desvio padrão, também observado na região Sudeste ( 16,8 ).
O menor Desvio padrão foi o da região Centro-Oeste ( 15,2 ).
Os demais desvios foram : 16,2 na região Norte, 16,5 nas regiões Nordeste e Sul.
A Mediana das notas dos estudantes de todo o Brasil foi 53,0.
A maior Mediana ocorreu na região Nordeste ( 57,9 ), e a menor, na região Sul ( 49,1 ).
As demais medianas foram : 51,5 na região Norte, 54,7 na região Sudeste e 57,6 na região Centro-Oeste.
A nota Máxima do Brasil, como um todo, foi 94,5, sendo obtida por, pelo menos, um aluno na região Nordeste.
As demais notas máximas foram : 85,1 na região Norte, 91,8 na região Sudeste, 90,3 na região Sul e 84,0 na região Centro-Oeste.
A nota Mínima foi zero na região Sudeste, e nas demais regiões, as notas mínimas foram : 3,7 na região Sul, 12,1 na região Norte, 14,1 na região Nordeste e 18,4 na região Centro-Oeste, a maior nota Mínima.
Observa -se que existe diferença estatisticamente significativa da menor Média das notas do Componente de Conhecimento Específico, obtida na região Sul ( 48,9 ), em relação às médias das regiões Nordeste, Sudeste e Centro-Oeste.
Assim como nos Gráficos 6.1 e 6.2, no Gráfico 6.3, apresentado a seguir, é possível uma avaliação do desempenho de concluintes de Artes Visuais ( Licenciatura ) em relação ao Componente de Conhecimento Específico com um histograma da distribuição das notas correspondentes.
Essa também é uma distribuição unimodal, e o grupo modal é o ( 50 ; 60 ], mesmo grupo modal da prova como um todo e de Formação Geral.
Na Tabela 6.6, é apresentada uma comparação dos resultados em relação à Categoria Administrativa, à Organização Acadêmica e a Modalidade de Ensino, agora levando em conta o desempenho de alunos do Componente de Conhecimento Específico da prova.
No que se refere à Organização Acadêmica, a maior Média foi a dos CEFET/IFET ( 56,4 ), vindo a seguir a das Universidades ( 54,3 ), depois, a das Faculdades ( 52,8 ) e, por fim, a dos Centros universitários ( 49,1 ).
O maior Desvio padrão, e acima do valor para o Brasil, como um todo ( 16,8 ), foi o dos CEFET/IFET ( 17,0 ) e, o menor desvio foi encontrado nas Universidades e Faculdades ( 16,4 ).
Alunos de Universidades obtiveram a maior nota Máxima ( 94,5 ).
Os Centros Universitários obtiveram nota Máxima igual a 91,4, as Faculdades obtiveram nota Máxima igual a 83,4, e nos CEFET/IFET, a nota Máxima foi 87,0.
As medianas foram : 48,6 nos Centros Universitários, a menor delas, 53,0 nas Faculdades, 55,8 nas Universidades e 59,9 nos CEFET/IFET, a maior.
A nota Mínima foi zero nos Centros Universitários, e a maior nota Mínima, 17,1, ocorreu nos CEFET/IFET.
Observa -se que existem diferenças estatisticamente significativas ao nível de 95 \% no Componente de Conhecimento Específico entre a menor média, obtida nos Centros Universitários ( 49,1 ), e as médias dos outros três tipos de Organização Acadêmica.
Quanto à Categoria Administrativa, observa -se um comportamento semelhante àquele da parte de Formação Geral e à prova como um todo, ou seja, existe diferença estatisticamente significativa entre as médias das IES Públicas ( 57,0 ) e IES Privadas ( 50,0 ).
Nesse caso, também, a maior Média foi obtida por alunos de IES Públicas de ensino.
Também semelhante é o comportamento das médias no que tange à Modalidade de Ensino, havendo à diferença estatisticamente significativa entre a Média da Educação Presencial ( 56,7 ) e a da Educação a Distância ( 47,0 ). 6.2 ANÁLISE DAS QUESTÕES OBJETIVAS Esta seção apresenta estatísticas selecionadas das Questões Objetivas de Formação Geral ( 6.2.1 ) e de Conhecimento Específico ( 6.2.2 ).
São também apresentadas e comparadas as médias das subpopulações caracterizadas por Grande Região. 6.2.1 Componente de Formação Geral Na Tabela 6.7, são apresentadas as Estatísticas Básicas relativas às oito questões objetivas do componente da prova que abrange a Formação Geral dos estudantes.
A Média do Brasil foi 45,5.
A menor Média foi encontrada na região Sul ( 43,9 ), e a maior, na região Centro-Oeste ( 50,9 ).
As demais médias foram : 44,7 na região Norte, 47,2 na região Nordeste e 46,1 na região Sudeste.
O Desvio padrão do Brasil foi 19,7, sendo o maior Desvio padrão encontrado na região Centro-Oeste ( 22,0 ), e o menor, na região Norte ( 18,8 ).
Os demais desvios foram : 20,1 na região Nordeste, 19,8 na região Sudeste e 19,1 na região Sul.
As medianas do Brasil, como um todo, e as das cinco regiões foram 50,0.
A nota Máxima 100,0 foi alcançada em quatro das Grandes Regiões.
A exceção foi a região Norte com nota Máxima 87,5.
As notas mínimas ( 0,0 ) foram iguais para todas as regiões.
Na Tabela 6.8, apresentam -se o Índice de Facilidade e o Índice de Discriminação ( ponto-bisserial ) para cada uma das questões objetivas do Componente de Formação Geral.
Quanto ao Índice de Facilidade, foram usadas as seguintes cores para diferenciar o nível de dificuldade da questão : • Azul para as questões classificadas com índice Muito fácil ( \&gt ; =0,86 ), verde para as questões classificadas com índice Fácil ( 0,61 a 0,85 ), amarelo para as questões classificadas com Médio ( 0,41 a 0,60 ), vermelho para as questões classificadas com Difícil ( 0,16 a 0,40 ) e roxo para as questões classificadas com Muito difícil ( \&lt ; =0,15 ).
Já quanto ao Índice de Discriminação, foram usadas as seguintes cores para qualificar a questão : • As questões classificadas com índice Fraco receberam a cor vermelha ( \&lt ; =0,19 ), as classificadas com Médio receberam a cor amarela ( 0,20 a 0,29 ), as classificadas com Bom receberam a cor verde ( 0,30 a 0,39 ) e as classificadas com Muito bom ( \&gt ; =0,40 ) receberam a cor azul.
As questões objetivas do Componente de Formação Geral, segundo o Índice de Facilidade, foram assim avaliadas : das oito questões, nenhuma teve o Índice de Facilidade classificado como Muito fácil, e duas questões foram consideradas como Fácil, por terem índice de acertos de 0,68 e 0,81 ( questões 4 e 6, respectivamente ).
Três questões foram consideradas com índice de dificuldade Médio, situando -se no intervalo entre 0,41 e 0,60 do Índice de Facilidade, ou seja, houve entre 41,0 \% e 60,0 \% de acertos, enquanto duas questões foram classificadas na categoria Difícil, situando -se no intervalo entre 0,16 e 0,40.
Por fim, uma questão apresentou menos de 15 \% de acertos, razão pela qual foi classificada como Muito difícil.
O Índice de Facilidade variou de 0,11 a 0,81, e o de Discriminação, de 0,25 a 0,53. facilidade, com uma proporção de 0,68 acertos.
O máximo de acertos foi alcançado pela questão 6, com um Índice de Facilidade de 0,81.
Duas questões tiveram Índice de Discriminação abaixo do Muito Bom : a questão 3 foi considerada com índice Bom, e a questão 1, com índice Médio.
Os Índices de Facilidade de tais questões foram : 0,25 para a questão 3 ( Difícil ) e 0,11 para a questão 1 ( Muito Difícil ).
Como, apesar de ser considerada Muito Difícil, a questão 1 teve Índice de Discriminação Médio, ela não foi eliminada pelo critério de ponto-bisserial do cálculo da nota de Formação Geral.
No Gráfico 6.4, para exemplificar, analisa -se o comportamento da questão de número 4 de Formação Geral.
Trata -se de uma questão considerada Fácil, com relação à facilidade e a que obteve o maior Índice de Discriminação dessa parte da prova.
Neste gráfico, cada uma das cinco curvas representa o percentual de respostas em determinada alternativa da questão, em função do número de acertos dos estudantes nessa parte da prova ( Formação Geral/Múltipla Escolha ), antes de possíveis eliminações pelo critério do ponto-bisserial.
A curva em verde corresponde à alternativa B, a alternativa correta para essa questão.
Observa -se que entre os estudantes com menor número de acertos, nessa parte do exame, a situação mais frequente foi a escolha da alternativa E ( em vermelho ), incorreta.
Por exemplo, entre os estudantes que acertaram três questões, mais da metade ( 63,8 \% ) escolheu a alternativa correta B ( em verde),15,5 \% escolheram a alternativa E ( vermelho ), 11,4 \% escolheram a alternativa D ( roxo ), 6,4 \% escolheram a alternativa A ( azul ) e 2,6 \% a C ( preto ).
Entre os que acertaram três respostas entre as questões de múltipla escolha de formação geral, 0,0 \% deixou a questão em branco e 0,4 \% marcou mais de uma alternativa, invalidando a questão.
À medida em que o número de acertos aumenta, indicando desempenho melhor nessa parte da prova, aumenta concomitantemente a proporção de estudantes que selecionaram a alternativa correta B, atingindo 100 \% para os estudantes com oito acertos.
Essa análise permite verificar como a questão discriminou os grupos de desempenho, justificando o alto índice obtido na questão.
Cumpre notar que não é possível inferir deste gráfico nem o índice de facilidade ( que seria uma média da proporção ponderada pela quantidade de alunos com cada uma das notas ), nem o índice de discriminação ponto-bisserial, por razão equivalente.
No caso extremo no qual a grande concentração dos acertos dos alunos fosse abaixo de 4, o índice de facilidade seria obrigatoriamente abaixo de 40 \% ( nesse exemplo ).
Caso a concentração fosse em 6 acertos ou mais, o índice seria obrigatoriamente acima de 60 \%.
Os gráficos relativos às demais questões de Formação Geral constam do Anexo I. 6.2.2 Componente de Conhecimento Específico Na Tabela 6.10, são apresentadas as Estatísticas Básicas em relação às questões objetivas do Componente de Conhecimento Específico da prova de Artes Visuais ( Licenciatura ) por Grande Região.
Nessa parte da prova, duas questões foram anuladas pela Comissão Assessora da Área, questões 16 e 22.
Além disso, como será discutido ainda nesta seção, duas questões objetivas do Componente de Conhecimentos Específicos não foram usadas no cômputo das notas, por terem sido descartadas pelo critério do ponto-bisserial.
Assim, as notas foram calculadas com base em 23 das 27 questões objetivas de Conhecimentos Específicos.
A Média do Brasil desse componente foi de 55,0.
A menor Média foi observada na região Sul ( 51,2 ), e a maior, na região Centro-Oeste ( 59,4 ).
O Desvio padrão de todo o Brasil foi 18,1, sendo o menor Desvio padrão encontrado na região Centro-Oeste ( 16,6 ), e o maior, na região Sudeste ( 18,1 ).
A Mediana de todo o Brasil foi 56,5, a mesma da região Sudeste.
Nas regiões Norte e Sul a Mediana foi menor que a nacional ( 54,4 e 52,2, respectivamente ).
As regiões Nordeste e Centro-Oeste registraram a maior Mediana ( 60,9 ).
A nota Máxima da prova foi 100,0, obtida nas questões objetivas do Componente de Conhecimento Específico, por, pelo menos, um aluno das regiões Nordeste e Sudeste.
A região Centro-Oeste registrou a menor nota Máxima, 87,0.
Pelo menos um estudante tirou a nota Mínima zero nas regiões Sudeste e Sul.
Nas demais regiões, a nota Mínima, em ordem crescente, foi : 8,7 na região Norte, 13,0 na região Nordeste e 21,7 na região Centro-Oeste.
Na Tabela 6.11, apresentam -se os Índices de Facilidade e Discriminação ( ponto-bisserial ) das questões objetivas do Componente de Conhecimento Específico para os estudantes de Artes Visuais ( Licenciatura ).
Para facilitar a diferenciação das questões, usaremos as mesmas cores da Tabela 6.8 para as diferentes classificações dos Índices de Facilidade e de Discriminação.
Observa -se que duas questões foram anuladas pela CAA, e que as análises dizem respeito às 25 questões válidas.
A partir do Índice de Facilidade obtido, pode -se concluir que 28 \% das questões objetivas da prova foram consideradas de nível Difícil, sete das 25 questões, foram consideradas válidas, e nenhuma foi classificada como Muito difícil.
Dez questões foram classificadas como Fácil, e outras oito consideradas como Médio.
Nenhuma teve mais do que 85 \% de acertos e, portanto, não houve questão classificada como Muito fácil.
Na Tabela 6.12, é da a distribuição das questões, considerando -se simultaneamente a classificação dos dois índices.
Dentre as questões que alcançaram os maiores índices de discriminação, classificadas com Muito bom, uma delas, a de número 18, foi a de índice mais Já quanto aos índices de discriminação das questões objetivas do Componente de Conhecimento Específico da prova, tem -se como resultado a seguinte classificação : 11 das 25 questões foram consideradas com índice Bom, enquanto nove delas tiveram Índice de Discriminação Muito bom.
Assim, para 80 \% das questões – 20 em 25 – os Índices de Discriminação foram Bom ou Muito bom.
Dentre as demais, três delas foram classificadas como Médio, e outras duas como Fraco, sendo cinco, por conseguinte, a quantidade de questões nos dois patamares mais baixos de discriminação.
Constata -se, assim, que a prova – no que se refere ao Componente de Conhecimento Específico – possuía boa capacidade de discriminar entre aqueles que dominam ou não o conteúdo.
Na Tabela 6.12, é da a distribuição das questões, considerando -se simultaneamente a classificação dos dois índices.
Dentre as questões que alcançaram os maiores índices de discriminação, classificadas com Muito bom, uma delas, a de número 18, foi a de índice mais elevado ( 0,50 ), tendo sido classificada como Fácil quanto ao Índice de Facilidade, com 66 \% dos estudantes marcando a opção correta.
A questão de número 28 foi a mais difícil dentre as 25 questões específicas válidas, com o mais baixo Índice de Facilidade, 23 \% de acertos.
Essa questão apresentou poder discriminatório muito baixo, 0,02, o que comprova ter sido a mais difícil para os estudantes.
Destaca -se, também, a questão 25, com Índice de Facilidade 0,25, o que, em termos percentuais, corresponde a 25 \% de estudantes que responderam acertadamente.
Já 0,18 foi o seu de Índice de Discriminação.
Tais questões foram, portanto, pelo critério ponto-bisserial, consideradas inadequadas.
Por isso, as questões 25 e 28 foram eliminadas do cômputo da nota final.
A título de exemplo das análises do comportamento das questões objetivas, o Gráfico 6.5 analisa -se a questão 18 do Componente de Conhecimento Específico.
Essa questão foi considerada pelas respostas dos estudantes avaliados na prova como uma questão com nível de facilidade Fácil, 0,66, ou seja, 66 \% dos estudantes assinalaram acertadamente a opção D, correspondente ao gabarito.
Como já comentado, seu Índice de Discriminação foi igual a 0,50, classificado como Muito bom, o maior valor de discriminação.
Nesse gráfico, cada uma das cinco curvas representa o percentual de respostas em determinada alternativa da questão 18, em função do número de acertos dos estudantes nessa parte da prova, antes de possíveis eliminações de questões pelo critério do ponto-bisserial.
A alternativa correta D, representada no gráfico pela curva em roxo, foi escolhida em maiores proporções pelos alunos com desempenho melhor nessa parte da prova.
Já as alternativas incorretas, também denominadas distratores, foram selecionadas, principalmente, por aqueles com notas mais baixas.
Observa -se que a soma não é 100 \% por causa das questões não respondidas ou com mais de uma opção marcada.
Aqueles com nota zero, na sua quase totalidade, deixaram essa questão em branco ou marcaram mais de uma alternativa, comportamento considerado inválido.
A proporção de alunos que selecionou a resposta correta D aumenta gradativamente, chegando a atingir 100 \% para 21 acertos ou mais, enquanto a proporção dos que escolheram alternativas incorretas decai, a partir de três ou quatro acertos, como função do número de acertos nessa parte da prova.
Os gráficos relativos às demais questões do Conhecimento Específico constam do Anexo I. 6.3 ANÁLISE DAS QUESTÕES DISCURSIVAS Esta seção apresenta estatísticas selecionadas e histogramas das Questões Discursivas de Formação Geral ( 6.3.1 ) e Conhecimento Específico ( 6.3.2 ).
São também apresentadas e comparadas as médias de subpopulações caracterizadas por Grande Região. 6.3.1 Componente de Formação Geral As análises dos resultados de desempenho dos estudantes de Artes Visuais ( Licenciatura ) nas duas questões discursivas relativas à Formação Geral encontram -se na Tabela 6.13 e no Gráfico 6.6.
Na Tabela 6.13, observa -se que a nota Média nesse conjunto de questões foi acima da obtida nas objetivas.
Os estudantes, de todo o Brasil, obtiveram, em Formação Geral, Média 45,5 nas questões objetivas e 54,3 nas questões discursivas.
Pode -se notar também que o Desvio padrão nesse conjunto de questões foi maior do que o obtido nas objetivas, 19,7 nas questões objetivas e 23,9 nas questões discursivas.
A maior Média foi obtida na região Norte ( 55,2 ), e a menor, na região Nordeste ( 51,9 ).
A Mediana de todo o Brasil, neste componente, foi 59,5.
Na região Centro-Oeste, a Mediana foi a menor ( 56,3 ), e a maior Mediana foi obtida na região Norte ( 61,5 ).
A nota Máxima foi 97,0, foi obtida por, pelo menos, um estudante da região Sudeste.
A menor nota Máxima, 94,0, foi obtida na região Norte.
A nota Mínima ( 0,0 ) foi a mesma em todas as regiões do Brasil.
No Gráfico 6.6, está representada a distribuição das notas nas questões discursivas do Componente de Formação Geral.
A moda dessa distribuição ocorre no intervalo ( 60 ; 70 ], seguida do intervalo ( 70 ; 80 ], o segundo com a maior frequência.
Destaca -se também o intervalo [ 0 ; 10 ] com frequência em torno de 8 \% do total de notas, sendo que no intervalo [ 0 ; 10 ] inclui -se, além da nota zero, a frequência de alunos que deixaram as respostas a esse tipo de questão em branco.
Na sequência, os resultados verificados para cada uma das questões discursivas de Formação Geral serão apresentados, estabelecendo -se relações com os conteúdos abordados em cada uma delas.
Os comentários da Banca de docentes corretores, a respeito do observado na correção das respostas dos estudantes, suas impressões e conclusões serão apresentados junto à análise de cada questão.
Cumpre esclarecer que, tendo em vista que as questões discursivas de Formação Geral são padronizadas, ou seja, constam de todas as provas, os comentários da Banca são os mesmos para todas as carreiras acadêmicas, sendo direcionados a todos os estudantes que participaram do Enade/2017.
A seguir, serão analisados os desempenhos da Área de Artes Visuais ( Licenciatura ) nas duas questões discursivas de Formação Geral do Enade/2017, comparando -se os resultados obtidos com comentários para cada questão. 6.3.1.1 Análise de Conteúdo da Questão Discursiva 1 do Componente de Formação Geral Os dados de Artes Visuais ( Licenciatura ), obtidos a partir das respostas à questão 1, encontram -se na Tabela 6.14 e no Gráfico 6.7.
Nas duas questões discursivas de Formação Geral, o desempenho dos alunos do Brasil como um todo, foi bem parecido : Média 53,8 na primeira e 54,0 na segunda.
A maior Média para a questão 1 foi obtida na região Norte ( 54,8 ), e a menor, na região Nordeste ( 51,2 ).
Quanto à variabilidade das notas, o Desvio padrão de todo o Brasil foi 28,1, menor do que o da questão 2.
O menor Desvio padrão foi obtido na região Sul ( 26,7 ), e o maior Desvio padrão foi obtido na região Nordeste ( 30,6 ).
As medianas do Brasil como um todo e as de três das regiões foram iguais ( 60,0 ).
Nas regiões Sul e Centro-Oeste as medianas foram menores, iguais a 55,0.
As notas máximas e mínimas da questão discursiva 1 foram, sem exceção, as mesmas para todas as regiões do Brasil, respectivamente, 100,0 e 0,0.
No Gráfico 6.7, mostra -se a distribuição das notas na questão discursiva 1 do Componente de Formação Geral.
Observa -se que a maior frequência corresponde aos alunos que obtiveram notas no intervalo ( 70 ; 80 ], com pouco mais de 20 \% dos participantes, seguido do intervalo ( 50 ; 60 ], com em torno de 17 \%.
O percentual dos que deixaram a resposta a essa questão em branco ficou em torno de 8 \% dos participantes, configurando -se como um máximo local. 6.3.1.2 Comentários sobre a correção de Conteúdo das respostas à Questão Discursiva 1 O enunciado era claro e indicava um tema de extrema relevância, qual seja, a incidência de doença grave integrante do rol das DST ( Doenças Sexualmente Transmissíveis ), que encaminha para a análise do direito fundamental à saúde, que tem \textbf{assento} constitucional ( CRFB, art. 6º. caput : São direitos sociais a educação, a saúde, a alimentação, o trabalho, a moradia, o transporte, o lazer, a segurança, a previdência social, a proteção à maternidade e à infância, a assistência aos desamparados, na forma desta Constituição ).
A questão focalizava, por meio de três textos, um tema social extremamente relevante, qual seja, a relação entre homem e mulher nas suas interações íntimas, no caso representadas pela vida sexual, e o reflexo das relações de poder veiculadas pela sociedade, que apontam, no mais das vezes, o predomínio do homem.
Tal posição de superioridade é justificada pelas convenções arraigadas que acarretam a educação dos homens e mulheres por padrões patriarcais, que redundam em posições machistas.
Um dos textos demandava uma reflexão sobre a realidade dos \textbf{casais} que convivem em relações estáveis ou fugazes e não têm o controle dos efeitos das relações sexuais, que abrangem a gravidez não planejada e, em decorrência do não cuidado, a inserção de doenças transmitidas sexualmente.
Esse panorama é refletido na pouca influência das mulheres sobre os homens quanto à proteção, normalmente representada pelo uso de preservativos nas relações intimas, e após a constatação da doença, no encaminhamento das pessoas para tratamento.
Como efeito colateral, está a transmissão de doenças para os fetos.
Outro texto trazia à tona o maior cuidado que as mulheres têm, em geral, com sua saúde, notadamente no momento em que se descobre grávida e que redunda na realização de exames obrigatórios de pré-natal, que podem facilmente detectar doenças e, caso encontradas, ter o seu tratamento encaminhado.
A par disso, destaca -se que a crise de recursos financeiros agravou a situação, que já recebe influxos sociais negativos, como a dificuldade do homem em utilizar preservativos, realizar exames e procurar tratamentos quando adoece.
Como muito divulgado pelas mídias, a crise financeira do país repercutiu diretamente na prestação de serviços de saúde em todas as esferas : distrital, federal, estadual e municipal.
Assim, agravou -se o problema da falta de medicamentos essenciais para o tratamento da sífilis e, certamente, de outras doenças.
Como se verifica, a temática indicada no enunciado, representada pelos textos, pertence à realidade conhecida amplamente pelos estudantes, pois também é divulgada por todos os meios de comunicação.
Trata -se de tema geral que, no entanto, normalmente, não se inclui nas grades curriculares de cursos superiores que não sejam da área de Saúde.
Apesar disso, a temática é de conhecimento geral e relacionada aos direitos humanos, esses sim, ponto de integração com as diretrizes curriculares.
A linguagem utilizada era clara, indicando a existência dos fatos vinculados a questões relevantes e à saúde pública.
No entanto, como os textos trazem uma ampla gama de possibilidades de interpretação, ocorreram algumas compreensões diversas das esperadas.
Além disso, ressalta -se que a amplitude dos textos disponibilizados no enunciado fornece elementos que poderiam ser simplesmente reproduzidos pelos concluintes nas suas respostas.
O enunciado da questão encaminha a resposta para dois aspectos.
O estudante deveria abordar conteúdo pertinente à vulnerabilidade das mulheres às DST e o papel social do homem em relação à prevenção dessas doenças.
Para responder a esses aspectos, um número expressivo de concluintes utilizou somente os elementos dos textos motivadores, sem nada acrescentar.
A resposta deveria conter também uma segunda parte, a indicação de duas ações voltadas para o público masculino, a serem adotadas no âmbito das políticas públicas de saúde ou de educação, para a redução do problema.
Verificou -se que o desempenho dos estudantes nessa questão foi de médio para bom.
Ao se analisarem as respostas, constata -se que muitos foram os que apresentaram os tópicos presentes no padrão de resposta, amplo o suficiente para albergar as respostas apresentadas pelos concluintes.
Tal resultado pode ser justificado, tendo em vista tratar -se de tema bastante difundido nas diversas mídias, por ter sido apresentado de forma ampla pelos textos motivadores presentes no enunciado.
Por outro lado, o debate no ambiente universitário sobre as relações de poder entre homens e mulheres tende a ser marcante pelo movimento de empoderamento feminino, liderado por organizações sociais com ampla atuação nos estabelecimentos de nível superior.
De forma recorrente, as relações patriarcais entre homens e mulheres foram vinculadas ao modelo machista vigente.
Essas condições, segundo os concluintes, explicam a incidência da doença e a situação de vulnerabilidade das mulheres, bem como a função social dos homens.
Outras respostas apontaram o poder público, nas áreas de saúde e/ou educação, como órgãos responsáveis pela solução do problema.
Para os que se fixaram na área da saúde, foram frequentes a indicação da necessidade de mais medicamentos, vacinas, educação sexual, médicos e postos de atendimentos.
Alguns, no entanto, apenas indicaram ações individuais para o homem, e outros, ações individuais para as mulheres.
Também figuraram respostas que apresentavam ações da igreja como solução, além de outras que faziam referências a campanhas no seio das famílias. 6.3.1.3 Análise de Conteúdo da Questão Discursiva 2 do Componente de Formação Geral Na Tabela 6.15, mostra -se que a Média do desempenho dos estudantes na questão discursiva 2 foi 54,0, foi pouco superior ao da questão discursiva 1 ( Média 53,8 ).
A região Norte foi aquela cuja Média, também nessa questão, foi maior ( 55,9 ), e a de menor Média foi a região Centro-Oeste ( 51,5 ).
Quanto à variabilidade das notas, o Desvio padrão de todo o Brasil foi 32,2.
O maior desvio nessa questão foi obtido na região Nordeste ( 33,8 ), enquanto o menor foi obtido na região Norte ( 31,0 ).
A Mediana de todo o Brasil foi 70,0, a mesma sendo obtida nas regiões Norte, Nordeste e Sudeste.
Nas regiões Sul e Centro-Oeste, as medianas foram 50,0, em ambas.
As notas máximas ( 100,0 ) e as notas mínimas ( 0,0 ) foram as mesmas em todas as regiões, sem exceção.
Nas regiões Sul e Centro-Oeste, as medianas foram 50,0, em ambas.
As notas máximas ( 100,0 ) e as notas mínimas ( 0,0 ) foram as mesmas em todas as regiões, sem exceção.
No Gráfico 6.8, mostra -se a distribuição das notas na questão discursiva 2 do Componente de Formação Geral.
Observa -se que a maior frequência corresponde aos alunos que obtiveram nota dentro do intervalo ( 70 ; 80 ].
Destacam -se, também, os alunos que deixaram a resposta à questão em branco e os alunos que obtiveram nota dentro do intervalo ( 90 ; 100 ], caracterizando -se como máximos locais.
Nota -se, ainda, que as notas ficaram mais espalhadas em comparação à questão discursiva de número 1.
Isso pode ser constatado também pela comparação do desvio padrão das notas da questão discursiva 2 ( 32,2 ) e das notas questão discursiva 1 ( 28,1 ). 6.3.1.4 Comentários sobre a correção de Conteúdo das respostas à Questão Discursiva 2 O enunciado era claro e abordava um tema de extrema relevância : a questão de gênero.
Trata -se de uma proposta polêmica, pois abrange a análise da inserção social de minoria vinculada à questão de gênero.
O enunciado solicitava que o concluinte discorresse sobre a importância do nome para as pessoas transgêneras e propusesse uma medida, no âmbito das políticas públicas, para facilitar o acesso dessas pessoas à cidadania.
A respeito da primeira parte, a importância do nome, esperava -se que o estudante abordasse as dificuldades que o uso do nome civil ( nome de registro, nome de batismo, nome da certidão de nascimento, por exemplo ) pode acarretar para as pessoas transgêneras, além de demonstrar conhecimento sobre a importância do nome social, direito já reconhecido através, principalmente, de decisões judiciais e atos administrativos, como decretos e portarias que dispõem sobre o uso do nome social.
Em atendimento à segunda parte do comando, o participante deveria propor uma política pública associada ao direito à cidadania das pessoas transgêneras.
Sobre esse aspecto, esperava -se a compreensão do significado de política pública, ou seja, propostas de âmbito governamental, e não, iniciativas de Organizações Sociais, igrejas ou grupos sociais de qualquer tipo.
O nascimento das pessoas, no Brasil, é objeto de inscrição em \textbf{cartórios} de Registro Civil das Pessoas Naturais, órgão responsável por estabelecer a certificação pública dos dados de todos os indivíduos que não somente nascem, mas casam, eventualmente se divorciam e vêm a óbito.
Daí o nome civil.
Para flexibilizar a correção, outras nomenclaturas foram consideradas, a saber : identificação de nascimento, nome de batismo, nome de registro, nome próprio, nome verdadeiro, que diferenciam o nome civil daquele que o indivíduo deseja utilizar para traduzir a sua opção pessoal de gênero, que se convencionou denominar nome social.
Quanto ao entendimento do termo transgênero, deve -se destacar que o movimento de mudança do gênero é contemporâneo e não está necessariamente vinculado à mudança de sexo.
Existem várias hipóteses em jogo.
Uma é a que envolve cirurgias de mudança de sexo ( a menina extirpa os seus seios e inclui órgão genital por intervenção cirúrgica ; o menino inclui seios e extirpa o seu órgão genital e inclui o órgão feminino também por intervenção cirúrgica ).
Há outro movimento que envolve apenas as vestimentas e o linguajar corporal, além dos relacionamentos ( as meninas assumem o trajar, o andar, os gostos dos meninos e mantêm relacionamentos afetivos com meninas ; os meninos assumem o trajar, o andar, os gostos das meninas e mantêm relacionamentos afetivos com meninos ).
Em relação à discussão da primeira solicitação do comando da questão, destaca -se que todos têm nome civil, que foi escolhido pelos seus pais ou responsáveis e, na falta deles, por autorização do Juiz de Direito, pautado no sexo detectado no nascimento.
Já o nome social é aquele que indica a preferência do indivíduo na sua identificação nas relações sociais e que destoa do seu gênero de registro civil.
Ressalte -se que o Brasil não possui ainda legislação específica sobre pessoas transgêneras.
Observa -se um relativo progresso no reconhecimento de direitos, através, principalmente, de decisões judiciais e atos administrativos, como decretos e portarias que dispõem sobre o uso do nome social.
Sem uma lei que defina os procedimentos da alteração dos documentos para pessoas transexuais, essa parcela da população LGBT ( lésbicas, gays, bissexuais, transexuais e transgêneros ) é obrigada a procurar na Justiça o reconhecimento de sua identidade, em processos que podem ser longos e que dependem do olhar de determinados, profissionais, como médicos, psicólogos, assistentes sociais, bem como do entendimento dos \textbf{advogados}, defensores públicos, juízes e do Judiciário, como um todo.
Atualmente, tramita na Comissão de Direitos Humanos da Câmara dos Deputados o Projeto de Lei João Nery ( 5002/2013 ), dos deputados Jean Willys ( PSOL-RJ ) e Erika Kokay ( PT-DF ), que determina que o reconhecimento da identidade de gênero seja um direito do cidadão.
O projeto recebeu o nome do primeiro trans-homem operado no Brasil.
Grande parte dos respondentes considerou a incongruência entre a imagem corporal, representada pelo gênero assumido pelos trans, e o nome que apresenta nos seus documentos como causa de sérios constrangimentos e embaraços e as dificuldades para que esses sujeitos acessem direitos básicos, como saúde, educação e empregabilidade em uma sociedade heteronormativa e preconceituosa.
Muitos demonstraram reconhecer que a identidade e expressão de gênero, através do nome social, pressupõe que sejam respeitadas a dignidade da pessoa humana, a intimidade, a liberdade e o direito à cidadania plena.
No entanto, também foram encontradas, em quantidades significativas, respostas que demonstravam a existência de uma interpretação social conservadora que não admite a mudança de sexo nem do nome social.
Existe um acentuado preconceito contra todos os movimentos que buscam identificação social diversa da civil.
Isso inclui os variados grupos : gays, travestis, transgêneros, dentre outros.
Já em relação ao segundo tópico a ser abordado pelo concluinte, é notório que as pessoas trans encontram -se destituídas de direitos no campo das políticas públicas, como saúde, habitação, emprego, educação, justiça, dentre outros.
Nesse sentido, uma grande diversidade de medidas, no âmbito das políticas públicas, poderia ser proposta.
Ações voltadas para a capacitação de profissionais da educação, saúde, segurança pública, dentre outras áreas, para melhor atender às demandas e especificidades desse público alvo seriam necessárias.
Ainda em relação à educação, seriam também necessárias ações que contribuíssem para diminuir a evasão, pois o ambiente escolar é difícil para uma pessoa transgênera, por conta das situações discriminatórias e preconceituosas vindas dos colegas de classe, e mesmo das/dos professores e/ou de outros profissionais da área, não podem ser mais um \textbf{complicador} nesse processo.
Ainda no que tange ao assunto, percebe -se que, ao serem excluídas dos bancos escolares, as oportunidades no mercado de trabalho formal praticamente inexistem.
Um outro aspecto que se evidencia no cotidiano desse segmento da sociedade diz respeito às situações de violência sofridas pelas pessoas transgêneras.
Nesse sentido, algumas propostas foram elencadas para minimizar o problema e combater a violência.
Políticas envolvendo melhor formação de funcionários das polícias civil e militar, por exemplo, somariam esforços na qualificação das abordagens policiais e dos atendimentos prestados pelas forças de segurança, pois o respeito à orientação sexual e à identidade e expressão de gênero é, antes de tudo, o respeito à dignidade, à cidadania e à própria democracia.
Nesse sentido, revela -se essencial que, ao sofrerem agressões, sejam elas físicas, psicológicas ou morais no espaço público ou domiciliar, as pessoas trans possam ser atendidas, com respeito inclusive ao nome social, nas delegacias comuns e, inclusive, nas Delegacias Especiais de Atendimento à Mulher ( DEAM ), pois esse é o órgão institucional que possivelmente detém as qualificações necessárias para o atendimento ao gênero feminino ao qual as travestis e mulheres trans efetivamente pertencem.
Também ocorre uma latente inacessibilidade das pessoas trans à assistência jurídica.
Nesse sentido, é necessário implementar projetos de capacitação para profissionais do poder judiciário, como um todo, de modo a prestarem a assistência jurídica gratuita na defesa intransigente dos direitos sociais, civis, trabalhistas e políticos dessa população.
Dessa maneira, a Justiça agiria de forma mais acessível e sensível para o segmento, que terá suas especificidades observadas a começar pelo devido respeito do direito ao uso do nome social.
Nesse âmbito, o mais comum foi propor a adoção de medidas punitivas para quem viola o direito à autodeterminação de gênero e a elaboração de leis que garantam a mudança do nome e assegurem outros direitos para as pessoas trans.
Propostas envolvendo a formação de profissionais, como as mencionadas aqui, e de conscientização da população foram as mais frequentes.
Muitas outras enfatizavam o direito ao uso respeitado do nome social e a garantia de alteração do registro civil, independentemente da cirurgia de redesignação sexual ou transgenitalização.
Muitos questionaram o fato de não bastar a autodeclaração no processo de alteração de nome, enfatizando que o nome social é uma questão de princípio de dignidade e cidadania.
Apesar de pouco mencionada nas respostas corrigidas, políticas envolvendo a promoção da saúde são muito necessárias.
Cabe aos gestores de saúde ações direcionadas, não somente à prevenção e tratamento do HIV/Aids, mas também políticas de saúde que contemplem as especificidades dessa população, tais como o Ministério da Saúde, especialmente o SUS ( Sistema Único de Saúde ), incluir, de fato, o atendimento para pessoas trans para acompanhamento da terapia hormonal, que muitas demandam, respeitando -se o direito de utilização do nome social.
De modo geral, são muitas as políticas públicas a que precisam ser implantadas.
Urge a necessidade de uma agenda mais compromissada por parte dos gestores do nosso país com essa comunidade, que ainda é muito discriminada.
Logo, tais políticas devem contemplar suas necessidades básicas como o direito de acesso aos estudos, à capacitação para o mercado de trabalho, à profissionalização, ao acesso a bens e serviços de qualidade em saúde, habitação, segurança, cultura e assistência social e, certamente, o respeito ao uso do nome social, o que também contribui para essa inclusão.
Assim, as opções apresentadas no padrão de resposta funcionaram como exemplos de políticas públicas possíveis, já que seria difícil apresentar uma lista suficientemente ampla para abranger a grande variedade de respostas encontradas.
O desempenho dos estudantes nessa questão foi médio.
As respostas, de modo geral, indicaram a necessidade de uma atualização por parte dos concluintes quanto a temas essenciais e, no caso da presente questão, a análise sobre aspectos fundamentais da questão de gênero.
Muitos, possivelmente motivados pelos textos apresentados no enunciado, restringiram -se a discorrer sobre preconceito de gênero ou mesmo sobre o preconceito de modo mais geral, sem responder ao que foi solicitado no comando da questão.
Outros apresentaram discursos religiosos ou ideológicos, também sem tratarem do que foi solicitado. 6.3.1.5 Análise de Língua Portuguesa das Questões Discursivas do Componente de Formação Geral Os dados de Artes Visuais ( Licenciatura ), obtidos a partir das respostas às questões discursivas do Componente de Formação Geral, no que tange à Língua Portuguesa, encontram -se na Tabela 6.16 e no Gráfico 6.9.
Nesse aspecto, os alunos, de todo o Brasil, obtiveram Média 56,5.
A maior Média, em relação à Língua Portuguesa foi obtida na região Sul ( 56,8 ), e a menor, na região Nordeste ( 54,4 ).
Quanto à variabilidade das notas, o Desvio padrão de todo o Brasil foi 20,8.
O menor Desvio padrão foi obtido na região Sul ( 19,6 ), e o maior foi obtido na região Nordeste ( 23,5 ).
A Mediana das notas de Língua Portuguesa foi 62,5 para o Brasil como um todo, a mesma obtida em quatro regiões.
A exceção foi a Mediana da região Sul, igual a 60,0.
A nota Máxima para todo o Brasil foi de 95,0, com, pelo menos, um aluno tirando essa nota na região Sudeste.
Na região Norte, a nota Máxima foi 82,5, enquanto nas regiões Nordeste e Centro-Oeste, foi 87,5 e na região Sul foi 90,0.
Já a nota Mínima foi zero em todas as regiões do país.
No Gráfico 6.9, mostra -se a distribuição das notas de Língua Portuguesa do Componente de Formação Geral.
Observa -se que a maior frequência ( mais do que 30 \% ) corresponde aos alunos que obtiveram nota no intervalo ( 60 ; 70 ].
Merece destaque, também, o percentual de alunos que deixaram as respostas para ambas as questões em branco, em torno de 7 \% do total, caracterizando -se como máximo local. 6.3.1.6 Comentários sobre a correção das respostas de Formação Geral com respeito à Língua Portuguesa Os enunciados apresentados em 2017 não explicitaram a exigência de elaboração de um texto “ dissertativo ”, como nos anos anteriores.
O participante deveria fazer uma exposição de seus conhecimentos sobre cada um dos assuntos e estruturar seus textos de acordo com as características do registro formal adequado à situação comunicativa – avaliação de conhecimentos.
Essa configuração determina exigências quanto aos seguintes aspectos : adequação da seleção vocabular, desenvolvimento do conteúdo, organização lógica das ideias, estruturação sintática dos períodos, utilização de procedimentos de encadeamento textual e de referenciação, obediência às exigências morfossintáticas próprias da modalidade escrita da norma-padrão, respeito às regras ortográficas e às regras de acentuação gráfica.
O padrão de resposta utilizado na avaliação das questões 1 e 2 considerou os aspectos relevantes ao bom desempenho linguístico como competências distintas, de modo a permitir um mapeamento detalhado do domínio dos recursos disponíveis na Língua Portuguesa para a comunicação escrita formal.
Com base nesse objetivo, foram avaliados os seguintes aspectos : ( a ) estruturação textual condizente com o gênero solicitado e o modo de organização textual expositivo adequado ao gênero – essa competência envolve :  estruturação sintática condizente com o padrão da modalidade escrita formal da língua portuguesa de modo a garantir a clareza necessária ;  distribuição do conteúdo do texto em parágrafos, de modo a garantir a sua organização temática ;  utilização de operadores discursivos que contribuam para a progressão temática do texto, estabelecendo relações lógicas entre as ideias apresentadas, tanto do ponto de vista intrafrasal, como do interfrasal ;  utilização de procedimentos de referenciação lexical e pronominal que permitam a retomada de referentes textuais ;  utilização de sinais de pontuação que contribuam para a organização lógica da frase e do texto ;  inteligibilidade relacionada ao atendimento das exigências de estruturação textual.
Espera -se, portanto, que o participante recorra a procedimentos linguístico-discursivos para organizar seu texto, permitindo o encadeamento lógico entre suas partes de forma a garantir a progressão e a coerência textuais.
Isso significa que os seguintes procedimentos foram considerados inadequados, de acordo com o padrão de resposta proposto :  elaboração de frases fragmentadas que comprometam a estrutura lógico-gramatical do texto ;  sequência justaposta de ideias sem encaixamentos sintáticos, reproduzindo hábitos da oralidade ;  elaboração de frase com apenas oração subordinada, sem oração principal ;  emprego equivocado do conector ( preposição, conjunção, pronome relativo, alguns advérbios e locuções adverbiais ) comprometendo a expressão da relação lógica entre duas ideias, com prejuízo da clareza do texto ;  emprego do pronome relativo sem a preposição, quando obrigatória ;  repetição ou substituição inadequada de palavras sem o emprego dos recursos oferecidos pela língua ( pronome, advérbio, artigo, sinônimo ) ;  emprego inadequado dos pronomes relativos “ cujo(a) ” e “ onde ” ;  utilização inadequada dos sinais de pontuação, comprometendo a clareza textual ;  inteligibilidade relacionada ao atendimento das exigências de estruturação textual. ( b ) respeito às convenções ortográficas da norma-padrão da Língua Portuguesa – essa competência envolve o domínio das regras de acentuação gráfica e da grafia padrão das palavras ( com ausência de abreviaturas próprias da linguagem da internet ), de acordo com as convenções estabelecidas pela legislação em vigor e consubstanciadas no Vocabulário Ortográfico da Língua Portuguesa, editado pela Academia Brasileira de Letras. ( Vale notar que, nessa edição de 2017, não foi aceita a legislação anterior, no caso das regras relativas ao uso do hífen e da acentuação gráfica ).
Espera -se que o participante : • grafe corretamente as palavras ; • respeite as regras de acentuação gráfica ; • empregue maiúsculas em início de frase, em nomes próprios de pessoas, lugares ou instituições ; • grafe as siglas com letras maiúsculas ou apenas com a primeira letra maiúscula quando formarem uma palavra ; • evite abreviações como p/, vc, tb, pq, tá, né, usadas muitas vezes em escrita informal e na internet ; • obedeça às regras de separação de sílabas no final da linha. ( c ) domínio dos diferentes aspectos morfossintáticos próprios da modalidade escrita formal da norma-padrão da Língua Portuguesa – essa competência envolve : a concordância nominal, a concordância verbal, a regência nominal, a regência verbal, a flexão nominal, a flexão verbal, a correlação entre os tempos verbais, a colocação pronominal e a utilização de sinais de pontuação que contribuam para a organização lógica da frase e do texto.
Espera -se que o participante : • flexione o verbo para estabelecer concordância de número com o sujeito da frase ; • flexione o artigo, o adjetivo e o pronome para concordar em número e em gênero com o substantivo a que se referem ; • observe a regência nominal e a verbal, utilizando a preposição adequada depois de um substantivo, um verbo ou um adjetivo ; • empregue adequadamente o acento grave indicador da crase entre uma preposição e um artigo ( a+a ) ; • obedeça às regras de colocação pronominal ( próclise e ênclise ), distintas dos hábitos da oralidade ou da escrita informal ( exigência de próclise com termo atrator, não exigência de que o pronome oblíquo se ligue ao verbo auxiliar por meio de hífen ) ;  flexione adequadamente verbos, substantivos, adjetivos e pronomes no que diz respeito à expressão das categorias gramaticais ;  flexione os verbos para expressar a correlação de modo e tempo nas estruturas subordinadas.
Com base nesses critérios, foram considerados como desvios de caráter morfossintático, e não como desvios ortográficos, as alterações que envolvem mudança de classe gramatical ou de forma flexional do verbo :  eliminação da marca de infinitivo ( -r- ) e substituição por acento agudo ou ausência total de marca do infinitivo ;  confusão entre “ ão ” e “ am ” nas formas verbais ;  confusão entre “ há ” e “ a ” ;  uso de hífen para separar pronome átono – tanto uso indevido quanto omissão ( exemplo : “ esperasse ”, em lugar de “ espera -se ” ; “ falar -mos ”, no lugar de “ falarmos ” ;  “ esta ” ( no lugar de “ está ” ) ; “ mais ” ( no lugar de “ mas ” ) ; “ e ” ( no lugar de “ é ” ).  verbos “ ter ” e “ vir ” que, na terceira pessoa do plural, não apresentarem o acento circunflexo, serão considerados inadequados como desvio de concordância, em aspectos morfossintáticos. ( d ) seleção vocabular adequada à modalidade escrita formal da Língua Portuguesa, exigida pela situação comunicativa – essa competência envolve a precisão na seleção/utilização do vocabulário relacionado à temática solicitada pela questão ; a ausência de marcas da oralidade, como termos de sentido muito genérico ( “ coisa ”, “ negócio ”, “ você ” ) e termos de registros mais informais ( como gírias, jargões, frases feitas, ditados populares, termos regionais ).
Assim, espera -se que o participante respeite a adequação vocabular não usando gírias ou expressões coloquiais, evite repetição desnecessária de palavras e utilize um vocabulário mais formal, como solicitado por um texto dissertativo. \textbf{Observações} : - A inteligibilidade relacionada ao atendimento das exigências de estruturação textual foi avaliada na segunda competência ( relativa aos aspectos textuais ). - Os problemas de coerência textual provocados por uso indevido do vocabulário foram avaliados na terceira competência ( relativa aos aspectos morfossintáticos e vocabulares ). - Não foram considerados como desvios morfossintáticos os problemas de caligrafia ( -a/ -o, -s, -r ). - Não foram considerados os textos grafados integralmente em caixa alta, já que não seria possível distinguir alguma marcação especial para as letras em início de frase. - Cada desvio foi considerado como uma ocorrência, mesmo que dois desvios fossem relativos ao mesmo aspecto linguístico.
A grade de avaliação do desempenho linguístico considerou, portanto, três grandes grupos de competências, segundo os aspectos explicitados anteriormente : a ) domínio das convenções ortográficas : grafia de vogais e consoantes, uso de maiúsculas e minúsculas, emprego do hífen e acentuação gráfica ; b ) domínio dos procedimentos de estruturação textual do ponto de vista microestrutural : organização interna dos períodos, emprego de conectores para a articulação lógica entre os períodos e entre os parágrafos, emprego de marcas de referenciação lexical e pronominal ; utilização dos sinais de pontuação que contribuem para a organização lógica da frase ; c ) domínio das regras de caráter morfossintático estabelecidas como modelares do ponto de vista da modalidade escrita formal da norma-padrão da Língua Portuguesa : concordância nominal e verbal, regência nominal e verbal, colocação pronominal, flexão nominal e verbal, correlação entre tempos e modos verbais, ausência de marcas de oralidade.
A seleção vocabular adequada à modalidade escrita formal da Língua Portuguesa foi incorporada a essa última competência, tendo em vista a intersecção entre as duas do ponto de vista das exigências do registro formal da modalidade escrita da norma-padrão.
Apreciação geral do desempenho dos estudantes : A correção revelou desempenhos distintos dos participantes, marcados pela falta de repertório cultural da maioria dos participantes : - a questão 1 propiciou a oportunidade de maior desenvolvimento do tema solicitado, gerando consequências na estruturação textual, já que os textos foram mais longos, mais elaborados e fluentes.
Entretanto, em virtude da diversidade de abordagens dos textos motivadores ( a falta de penicilina, a recusa no uso da camisinha pelos homens, a maior frequência das mulheres na realização de exames ginecológicos, a transmissão das DST aos bebês durante a gestação ), observaram -se respostas que se configuraram como verdadeiras paráfrases, com comprometimento do caráter \textbf{autoral} desejável em uma questão discursiva de formação geral.
Alguns participantes selecionaram um trecho de cada texto motivador para compor suas respostas, outros copiaram trechos inteiros. - a questão 2 teve um comportamento distinto, devido ao caráter polêmico que o caracteriza.
Assim, os textos são mais curtos, com vocabulário repetitivo e preso aos textos motivadores.
Além disso, observou -se grande ocorrência de protesto em função do tema, ora por motivos religiosos ora por preconceito de gênero.
Quanto aos aspectos linguísticos analisados durante esta avaliação, os resultados observados estão descritos nos parágrafos seguintes.
Aspectos ortográficos : O desempenho dos participantes revelou uma diferença muito grande nos dois aspectos analisados nesta competência : baixo índice de desvios da grafia padrão e grande índice de desvios de acentuação.
Em vários casos, ocorre ausência completa de acentuação gráfica.
Os resultados revelam que a tendência dominante entre os universitários brasileiros é a eliminação da acentuação gráfica, talvez motivada pelos hábitos relacionados às redes sociais e pela ausência de esclarecimento dos meios de comunicação, das autoridades e das escolas sobre as decisões do Acordo Ortográfico da Língua Portuguesa de 1990.
Os casos mais sistemáticos de eliminação do acento indicador da sílaba tônica são :  palavras proparoxítonas ( “ sifilis ”, “ proximos ”, “ politicas ”, “ publicas ”, “ transgeneros ” ) ;  palavras paroxítonas terminadas em ditongo crescente ( “ necessario ”, “ noticiarios ”, “ individuo ”, “ dependencia ”, “ varios ”, “ propria ”, “ transmissiveis ”, “ ocorrencia ” ) ;  palavras paroxítonas com hiato ( “ saude ” )  palavras oxítonas ( “ ninguem ”, “ esta ”, “ ate ”, “ tambem ”, “ prevencao ” ).
Por outro lado, destaca -se o uso indevido do acento gráfico em palavras como “ genêro ”, “ melâncolia ”, “ prevênção ”, “ intervênção ”, por exemplo.
Há, também, uma tendência a acentuar paroxítonos como se fossem oxítonos : “ transmissivéis ”.
Quanto ao domínio das convenções relativas à grafia das palavras, observam -se desvios recorrentes, como a hipercorreção pela escolha de “ e ” no lugar de “ i ”, por influência de hábitos da oralidade ( “ descriminação ” no lugar de “ discriminação ”, “ entervenção ” no lugar de “ intervenção ” ) ; desvios de grafia relacionados à variação diastrática podem ser observados em “ estrupo ”, “ subjulgadas ”, “ precoseito ”, “ soubre ”, “ vecendo ”, “ indesencia ”, “ apolojia ”, “ fulga ”, “ dereitos ”, “ sifelis ”.
Observam -se, também, casos de inadequação no uso da maiúscula para destacar determinadas palavras-chave do texto, como “ Transgêneros ”, “ Brasileiros ”, “ Homens ”.
Destaque -se, também, o grande número de participantes que grafam os textos inteiramente em caixa alta.
Vale observar, também, que, ao contrário do que se esperava, não apareceram abreviaturas próprias do “ internetês ”, relacionadas ao uso de redes sociais e e-mails.
Aspectos textuais : Esta competência é a que se revela como a mais problemática entre os participantes, porque são muitos os problemas observados, desvios acumulados durante toda a formação do estudante e que não se resolvem com um estudo autodidata, como acontece com regras ortográficas ou morfossintáticas.
São eles : sequência justaposta de ideias sem encaixamentos sintáticos ; redução drástica de estruturas subordinadas, ao lado do aumento na frequência de estruturas coordenadas e absolutas ; redução no uso de conectores para expressar relações lógicas essenciais à construção do texto, substituídas pela exigência de inferência por parte do interlocutor para suprir a sua ausência ; emprego equivocado de operadores que não estabelecem relações lógicas coerentes entre ideias do texto ; emprego inadequado do pronome relativo ( com omissão da preposição ou a utilização de pronome inadequado, como “ onde ” ) ; repetição exaustiva de termos sem a utilização de procedimentos mais sofisticados de substituição ( hiperonímias, hiponímias, nominalizações, expressões metafóricas ) ; frases fragmentadas que comprometem a estrutura lógico-gramatical ; frases formadas apenas por oração subordinada, sem oração principal.
Um importante aspecto a destacar é o baixo desempenho de uma parte dos participantes em relação à estrutura formal do texto produzido, o que é preocupante ao se levar em conta que são graduandos em fase final de formação.
São frequentes os casos de desvios de estruturação frasal, com uso inadequado ou ausência de conectivos entre parágrafos e entre frases.
Em uma parte dos textos, falta textualidade e domínio do registro padrão da língua.
Na verdade, observam -se relações linguísticas quase agramaticais, como as estabelecidas pela sequência de gerúndios sem o apoio de um ponto de partida para a organização das informações gramaticais e semânticas.
Observou -se que uma grande parte dos participantes não distribuiu as ideias em parágrafos, talvez devido ao pequeno número de linhas disponibilizadas para a resposta da questão ou, quem sabe, pela suposição de que não seria necessária essa divisão por não se tratar de um texto no modelo de uma redação dissertativo-argumentativa, como solicitado nos vestibulares.
Em função do encaminhamento dos enunciados das duas questões, que solicitaram uma análise do problema e encaminhamentos de políticas públicas, houve uma grande tendência, também, de construção de dois parágrafos desconexos, sem utilização de elementos coesivos adequados à progressão textual. solicitaram uma análise do problema e encaminhamentos de políticas públicas, houve uma grande tendência, também, de construção de dois parágrafos desconexos, sem utilização de elementos coesivos adequados à progressão textual.
Quanto à utilização dos mecanismos de referenciação, deve -se destacar a ocorrência de repetições de palavras ou expressões sem a utilização de termos sinônimos ou pronomes, como seria adequado.
Quanto à utilização dos sinais de pontuação, observou -se uma precariedade nos textos analisados.
É muito frequente a ocorrência de parágrafos sem marca interna de pontuação para separar os períodos.
Vale observar que não foi penalizada a ausência de vírgula para destacar locuções ou adjuntos adverbiais de pequena extensão deslocados de posição na frase, por ser um uso opcional.
São os seguintes os tipos de problemas encontrados : a ) vírgula : utilização de vírgula para separar o sujeito e o predicado ; ocorrência de apenas uma das vírgulas para separar uma palavra, uma expressão ou uma oração encaixada ; uso de vírgula no lugar do ponto para separar ideias que constituem períodos distintos ; ausência de vírgula para separar elementos de uma enumeração ; ausência de vírgula para separar oração adjetiva explicativa ou utilização inadequada para separar oração adjetiva restritiva ; b ) ponto e vírgula : utilização do ponto e vírgula no lugar de vírgula ; c ) ponto final : ausência de ponto final para separar períodos.
Aspectos morfossintáticos e vocabulares : Em relação à regência, o desvio mais frequente é a falta do sinal indicativo da crase – isso revela que o usuário não tem consciência de que, sob a forma do termo “ a ”, existe a presença de uma contração entre a preposição “ a ” ( exigida pela regência do termo anterior ) e o artigo definido “ a ”.
Um desvio de regência significativo, nos últimos anos, é a utilização inadequada de uma preposição ou sua ausência após o verbo ou o nome ( substantivo ou adjetivo ).
Outro problema relacionado à regência verbal e à nominal, encontrado frequentemente nas questões, foi a ausência de preposição antes de pronome relativo, processo generalizado na modalidade oral da língua, em situações de registro informal.
Apesar da possibilidade de que essa alteração de regência se generalize no padrão escrito da Língua Portuguesa, como já está ocorrendo até em textos jornalísticos, o não emprego da preposição foi considerado inadequado neste processo de avaliação.
A concordância verbal e a concordância nominal apresentam alguns desvios muito frequentes.
Quanto à concordância de número, observou -se ausência de marca ( com sujeito anteposto ou posposto ) ou uso indevido ( uso inadequado da marca de plural comandado pelo núcleo plural da locução adjetiva, apesar de o substantivo que funciona como núcleo do sintagma nominal estar no singular ).
Uma ocorrência que se destacou foi a ausência de acento circunflexo na forma plural do presente do indicativo do verbo “ ter ”, que foi considerada como um desvio na concordância verbal e não na acentuação gráfica.
Quanto à concordância de gênero, vários casos foram observados, normalmente no âmbito de sintagmas nominais longos, em que o adjetivo está afastado do substantivo.
Deve -se destacar o aparecimento da marca de plural em verbos ou adjetivos comandados por adjunto adnominal plural, apesar de serem relacionados a núcleos substantivos no singular, evidenciando um processo de hipercorreção ( exemplo : “ O nome social dos trangêneros podem ser estabelecidos por uma legislação específica ” ).
Quanto à questão da colocação pronominal, foram poucos os desvios observados.
Concluiu -se que, no registro escrito formal, a maioria dos participantes já incorporou regras como a não introdução da frase por um pronome oblíquo e a próclise na presença de um termo atrator.
Não se adotou, entretanto, o padrão excessivamente formal descrito pelas gramáticas normativas em relação à posição do pronome oblíquo em locuções verbais, já que esse uso está muito distante da prática cotidiana, até em textos mais formais.
Quanto aos aspectos vocabulares, alguns tipos de inadequação foram observados : expressões da oralidade ; seleção vocabular incompatível com o contexto, gerando falta de inteligibilidade ; falta de domínio de vocabulário mais abstrato e de maior complexidade, essencial ao desenvolvimento do texto de base dissertativa.
O principal aspecto observado foi a excessiva repetição de certas palavras, revelando limitação de repertório vocabular.
O termo “ pessoa ”, por exemplo, chega a ser repetido até 7 ou 8 vezes em um mesmo texto.
Várias marcas de oralidade foram identificadas, embora não com alta frequência : uso do pronome relativo “ onde ” como relativo universal, falta de artigo definido antes de substantivo, reduções como “ tá ”, “ pra ”, “ pro ”, “ prum ”, expressões informais, eliminação de preposições.
Em função do tema solicitado na questão 1, a sigla DST foi grafada de diferentes maneiras : DST´s, DSTs ou DSTS.
Além disso, houve flutuação de gênero gramatical nos sintagmas que envolviam a sigla : “ os diferentes DSTs ” ou “ as diferentes DSTs ”.
Em função do tema solicitado na questão 2, houve flutuação de gênero gramatical nos sintagmas que envolviam a sua palavra-chave : “ os trans ” ou “ as trans ” ; “ os transgêneros ” ou “ as transgêneras ”.
Todas essas formas foram aceitas, por considerarmos que sua utilização 6.3.2 Componente de Conhecimento Específico Na parte da prova relativa às questões discursivas do Componente de Conhecimento Específico ( Tabela 6.17 ), observa -se que a Média foi mais baixa do que para as questões discursivas do Componente de Formação Geral.
Enquanto no Componente de Formação Geral a Média para estudantes de Artes Visuais ( Licenciatura ) de todo o Brasil foi 54,3, na parte de Conhecimento Específico a Média foi 36,2.
A maior Média deste componente foi obtida pelos estudantes da região Nordeste ( 39,1 ), e a menor, pelos da região Centro-Oeste ( 34,7 ).
Quanto à variabilidade das notas, o Desvio padrão de todo o Brasil foi 19,7.
O maior Desvio padrão foi encontrado na região Nordeste ( 22,1 ), e o menor, na região Sul ( 18,7 ).
A maior nota Máxima, 93,3, foi obtida na região Nordeste.
Nas demais regiões a nota Máxima foi : 80,0 na região Norte, 90,0 na região Sudeste, 86,7 na região Sul e 73,3 na região Centro-Oeste.
A nota Mínima ( 0,0 ) foi obtida por pelo menos um aluno em todas as regiões do Brasil.
A Mediana do Brasil como um todo foi 36,7.
A maior Mediana foi obtida na região Nordeste ( 40,0 ), e a menor, na região Sul ( 33,3 ).
No Gráfico 6.10, representa -se a distribuição das notas nas questões discursivas no Componente de Conhecimento Específico.
A moda desta distribuição ocorre no intervalo ( 30 ; 40 ], com frequência em torno de 22 \% do total de participantes.
Observa -se que o intervalo [ 0 ; 10 ] é um máximo local, com frequência em torno de 12 \%.
Nesse intervalo estão computadas as ocorrências de respostas em branco para as três questões discursivas de Conhecimentos Específicos.
A análise de cada uma destas questões será feita a seguir. 6.3.2.1 Análise da Questão Discursiva 3 do Componente de Conhecimento Específico Na questão 3, cujos resultados aferidos encontram -se descritos na Tabela 6.18, a Média dos estudantes de todo o Brasil foi 26,6, a questão de mais baixo desempenho dentre as três discursivas de Conhecimentos Específicos.
A menor Média nessa questão foi obtida pelos alunos da região Centro-Oeste ( 18,7 ), enquanto a maior Média foi obtida na região Nordeste ( 37,0 ).
Quanto à variabilidade das notas, o Desvio padrão de todo o Brasil foi 32,5.
O maior Desvio padrão foi obtido na região Nordeste ( 35,9 ), enquanto o menor foi obtido na região Centro-Oeste ( 29,1 ).
A nota Máxima, 100,0 pontos, foi alcançada por, pelo menos, um aluno de quatro regiões.
A exceção foi a região Centro-Oeste, com nota Máxima 90,0.
A Mediana do Brasil como um todo foi 10,0, a mesma das regiões Norte, Sudeste e Sul.
A Mediana da região Nordeste foi a maior ( 30,0 ).
Já a da região Centro-Oeste foi zero, indicando que mais de 50 \% dos estudantes da região tiraram nota zero.
A nota Mínima ( 0,0 ) foi a mesma em todas as regiões do Brasil.
No Gráfico 6.11, mostra -se a distribuição das notas na questão discursiva 3, do Componente de Conhecimento Específico da área de Artes Visuais ( Licenciatura ).
A moda da distribuição foi a categoria dos estudantes deixaram a resposta à questão em branco, com um pouco mais do que 25 \% dos participantes.
Considerando que cerca de 20 \% receberam nota zero, quase metade dos estudantes não pontuaram nessa questão. 6.3.2.2 Comentários sobre as respostas à Questão Discursiva 3 Os conhecimentos exigidos pelo enunciado estão de acordo com o esperado, e o nível de profundidade extremamente razoável, já que não exigia uma reflexão profunda ou uma abordagem crítica e sim uma citação de um conteúdo de informação que é básico no ensino da Educação em Artes Visuais.
Para isso o enunciado foi claro na pergunta e objetivo na resposta esperada.
No entanto, o enunciado cita algumas referências que serviram de base para a elaboração da Abordagem Triangular e isso provocou um número significativo de desvio de respostas que não eram esperadas, respostas equivocadas, de quem optou por citar referências do próprio enunciado.
Por isso, os itens do enunciado não corroboraram para melhorar o aproveitamento das respostas corretas, mas sim atrapalharam.
O padrão de resposta foi adequado em relação ao enunciado, que pede para que sejam listadas e comentadas as três abordagens metodológicas.
A questão é bem objetiva com pouca ou quase nenhuma possibilidade de uma variação subjetiva de argumentação.
As possíveis variações se referem à utilização de termos com significados similares ( como ver, observar, ler, olhar, apreciar ) que tiveram espaço para serem desenvolvidos, para esclarecer o contexto da aplicação dos mesmos.
Em muitos casos, o termo específico não estava correto, mas sua função e/ou aplicação estavam corretamente compreendidas e explicadas.
Também foi possível registrar uma boa porcentagem de respostas com uma linguagem adequada aos termos específicos da área.
O padrão tem a clareza de uma resposta objetiva e oferece espaço para comentários que podem complementar e justificar possíveis interpretações.
A abordagem mais recorrente para responder à questão foi listar apenas as três ações da Abordagem Triangular sem os comentários solicitados, o que reduziu a média das notas.
Os dois principais erros foram : citar as referências do enunciado para responder à questão e listar os termos aprender - ensinar – educar como os pontos da metodologia.
Não houve mais diversidade nas respostas apresentadas e o desempenho dos estudantes foi insuficiente.
As respostas demostram que existe ainda uma falta de uniformidade na formação de professores na área.
Fica muito claro quando o estudante conhece o tema abordado ou não.
Em alguns casos de respostas incompletas, o estudante demonstra conhecer a autora A.
M.
Barbosa, mas afirma não conhecer a teoria em questão.
O enunciado faz referência a um conteúdo básico que deveria ser amplamente conhecido. 6.3.2.3 Análise da Questão Discursiva 4 do Componente de Conhecimento Específico Na Tabela 6.19, contam as informações relativas à questão 4 do conjunto de questões do Componente de Conhecimento Específico.
O desempenho dos estudantes nessa questão foi o melhor dentre as questões discursivas de Conhecimentos Específicos.
A Média geral do Brasil foi 49,2, sendo a menor Média registrada na região Nordeste ( 48,1 ), e a maior, na região Norte ( 50,2 ).
A nota Máxima ( 100,0 ) foi atingida por, pelo menos, um concluinte de todas as regiões.
A Mediana em todo o Brasil foi 60,0, sendo que o mesmo valor foi obtido nas cinco Grandes Regiões.
E, por fim, a nota Mínima de todas as regiões foi zero.
No Gráfico 6.12, está representada a distribuição de notas da questão discursiva 4, do Componente de Conhecimento Específico.
Essa distribuição tem moda no intervalo ( 50 ; 60 ] e um percentual de respostas em branco bem inferior ao da questão 3 e um pouco inferior ao da questão 5. 6.3.2.4 Comentários sobre as respostas à Questão Discursiva 4 A questão aborda conhecimentos básicos exigidos e apresenta nível de profundidade adequados ao contexto das diretrizes da área.
De modo geral, o enunciado utilizava linguagem clara e adequada ao nível de profundidade exigido, focando o papel do professor de Artes Visuais no ambiente escolar ; no fazer artístico articulado com produção de visualidades ; na diversidade de identidades ; na riqueza cultural ; e na pluralidade de valores.
A questão era composta por uma citação clara e adequada, mas o termo “ pós-crítico “ não parece ter sido tão claro e, de forma generalizada, não foi mencionado nas respostas.
O tema se abria para questões sobre objetos de conhecimento escolar ; produções poéticas ; diversidades ; cultura visual ; identidades juvenis ; e, questões pertinentes à contemporaneidade.
O enunciado continha também duas imagens fotográficas, mas as referências às imagens eram bem pequenas.
As respostas foram bem desenvolvidas, e a questão parece ter sido fácil para grande número de participantes.
O padrão de respostas ficou dentro das expectativas e foi possível contemplar boa porcentagem da variedade de respostas, mesmo que estas estivessem aplicando parte do enunciado nas soluções.
A abordagem mais frequente nas respostas foi a necessidade de aceitação da diversidade, considerando o professor como o mediador desse processo de aceitação.
Em segundo, a necessidade e a importância da diversidade para o crescimento e enriquecimento do processo de ensino-aprendizagem.
De maneira geral, as respostas repetem parte do enunciado, parafraseando -o sem elaborar os pontos citados no padrão de resposta.
Não houve muita diversidade nas respostas, mas sim uma forma dispersa de argumentar, sem objetividade.
O desempenho dos estudantes foi médio, o melhor da prova específica.
Percebe -se que a diversidade cultural e a democratização do ensino são temas recorrentes e bastante discutidos no ambiente escolar.
Também é notável a aprovação dessa necessidade com fator positivo e necessário.
No entanto, também é notável a dificuldade de objetividade na linguagem, o que, na maioria dos casos prejudicou a nota, pois se tornou impossível entender o que se desejava expressar. 6.3.2.5 Análise da Questão Discursiva 5 do Componente de Conhecimento Específico Na Tabela 6.20, constam as informações relativas à questão 5 do conjunto do Componente de Conhecimento Específico.
O desempenho dos estudantes de todo o Brasil nesta questão foi intermediário, mais elevado do que o da questão 3 e mais baixo do que o da questão 4.
A nota Média dos estudantes de todo o Brasil foi 33,0.
A maior Média foi registrada na região Centro-Oeste ( 35,7 ), enquanto a menor Média foi registrada na região Nordeste ( 32,3 ).
Quanto à variabilidade das notas, o Desvio padrão dos alunos do Brasil, como um todo, foi 20,5.
Enquanto o maior desvio foi encontrado na região Nordeste ( 23,2 ), o menor foi encontrado na região Sul ( 18,6 ).
A Mediana para o Brasil foi 30,0, a mesma de quatro regiões.
Apenas a região Centro-Oeste teve valor maior para a Mediana, igual a 40,0.
Para o conjunto de alunos de Artes Visuais ( Licenciatura ) do Brasil, a nota Máxima foi 100,0, obtida por pelo menos um estudante das regiões Norte e Sudeste.
A nota Mínima ( 0,0 ) foi a mesma em todas as regiões.
No Gráfico 6.13, representa -se a distribuição das notas da questão discursiva 5 do Componente de Conhecimento Específico.
Essa distribuição tem moda no intervalo ( 20 ; 30 ] e um máximo local na categoria “ em branco ”. 6.3.2.6 Comentários sobre as respostas à Questão Discursiva 5 A questão atendia a aspectos das Diretrizes Curriculares da Área de Licenciatura em Educação em Artes Visuais, mas partia de um nível complexo de análise apresentada no enunciado.
O texto tratava da abertura proporcionada pelas tecnologias da informação que injetam potencial criativo a um paradigma iniciado no contexto pós-industrial, no qual se constata que a mudança do valor de produção passar para a informação.
Tal paradigma alimenta -se ainda por uma revisão da condição do observador passivo para a de interativo, permitido por meio de obras assistidas por tecnologias e que apresentam comportamento, oferecendo uma base de conteúdo para se especular sobre a telemática como potencializador de obras abertas.
Ou seja, o texto discorria sobre uma mudança no estatuto da obra de arte, que se desloca da condição de obra fechada para obra aberta, participativa, interativa, na qual o significado não se encontra a priori no objeto de arte, mas sim na relação estabelecida com o observador e pelo observador ; daí, advém, como informado no texto, seu poder “ cognitivo ”, cultural e social.
Com pouquíssimas exceções, esse entendimento foi observado e apresentado nas respostas.
Talvez devido ao fato de o enunciado deslocar a atenção do estudante para questões mais específicas, como a “ importância do uso de dispositivos móveis no ensino de Arte ”, ou de apontar para questões sobre a “ democratização ”, possibilitada pelo uso de dispositivos móveis, no acesso à produção artística e sua circulação, a grande maioria das respostas discorreu sobre questões de acesso a museus, como sendo esse o grande benefício das redes, sem mencionar nada sobre o potencial das redes como elemento de constituição das obras em sua vocação comunicativa.
Foram respostas de cunho metodológico, abordando a tecnologia como ferramental didático, sem contemplar os conceitos acima descritos na relação entre Arte e Tecnologia.
O comando da questão solicitava que o texto a ser redigido pelo estudante abordasse três aspectos : a importância do uso de dispositivos móveis no ensino da arte ; a democratização no acesso à produção artística ; e as relações entre arte e tecnologia na modernidade.
Apesar de o termo “ modernidade ”, que nos parece equivocado, constar do enunciado, no padrão de resposta, usou -se o temo “ contemporâneo ”, o que nos parece mais adequado.
A maioria das respostas se restringiu ao enfoque da democratização, tanto na divulgação da produção dos respondentes, em sala de aula, quanto no acesso à produção artística, em geral, sem que fossem abordadas questões conceituais relacionadas à desconstrução de cânones e estereótipos da arte, ou à interação a nível criativo e constitutivo da obra, não apenas como dispositivo de acessibilidade ( museu virtual ou interação com o próprio aparelho celular ).
De modo geral, as respostas atenderam a um ou dois dos aspectos principais do padrão.
Dentre os desvios de padrão, mesmo com pouca incidência, os seguintes comentários não estavam previstos : aspectos punitivos do uso inadequado das tecnologias, em relação a direitos de imagem ; privacidade ; difamação ; jogos eletrônicos e suas contribuições para a inclusão social ; familiarização com os dispositivos ; divulgação de conteúdo ; plataformas de ensino em rede que utilizam dispositivos móveis para acesso e compartilhamento ; e ; ubiquidade da tecnologia e dispositivos móveis.
A média baixa se deu, possivelmente, em função da complexidade do enunciado.
Supõe -se que, por essa razão as respostas tenham apresentado, em sua maioria, uma cópia de palavras apontadas no texto guia e uma forte indicação do entendimento das questões sobre arte e tecnologia aquém da profundidade do tema.
Grande parte das \textbf{repostas} foram, no entanto, reveladoras sobre o estado da arte do ensino nesse campo, com o uso de mídias As poucas respostas que mereceram notas mais elevadas apresentaram um conhecimento mais amplo e sintonizado com as diversas manifestações artísticas contemporâneas que lidam com a característica policontextual e polissêmica de obras de nosso tempo.
O desempenho dos estudantes foi insuficiente.
Reiterando, ficou evidente que ocorreu uma incompleta compreensão da área de Arte e Tecnologia e de suas particularidades, talvez devido a ser um tema em ebulição, cuja história está ainda sendo escrita.
A tecnologia móvel é apenas uma faceta do tema, apesar da enorme relevância do assunto, motivo de festivais de arte sobre essa prática no mundo todo.
Considera -se, portanto, de suma importância a questão, não apenas por tratar de conteúdo contemporâneo, mas também por permitir tornar evidente o quanto o ensino na área ainda se ressente de investimentos na capacitação e atualização/formação de seus profissionais.
Constatou -se que as discussões sobre a prática e a teoria em Arte e Tecnologia ainda não fazem parte do universo acadêmico em geral.
Mas com a disseminação e a apropriação popular das tecnologias computacionais móveis e conectadas na internet, essa discussão torna -se a cada dia mais premente. 6.3.3 Considerações Finais Considerou -se a prova específica em Licenciatura em Artes Visuais do Enade 2017 uma avaliação que se propôs a cobrir aspectos contemporâneos e atuantes na área, e que, de modo geral, as questões foram elaboradas de forma a provocar respostas críticas e reflexivas.
O retorno foi positivo, com contribuições que apontam para uma boa diversidade na formação dos estudantes.
Os três temas, objeto das três questões discursivas, são importantes e formam um conjunto de posturas de grande relevância na atuação do professor hoje.
É interessante observar as médias muito próximas entre as Questões 3 e 5, sendo a primeira de caráter mais teórico e metodológico, e a segunda, focada na aplicação de tecnologias contemporâneas no ensino das Artes.
Era de se supor que a Questão 3 fosse de domínio geral, o que não prevaleceu, e, por outro lado, apesar das dificuldades específicas, a ubiquidade das tecnologias se fez presente e atuante.
A dificuldade de escrita por parte dos estudantes foi recorrente nas três questões discursivas da área.

% matched lemmas: advogado, assento, autoral, cartório, casal, complicador, observação, repor
\end{textsample}
